\section{Localization and the counting theorem}
\label{sec:localization}

In this section $M$ is closed, oriented and negatively curved,
$N=2n$, and $g$ is a regular negatively curved metric.
We prove the localization formula announced in
Section 1 and the quantitative counting theorem.

We first calculate the local degree at a point of maximal multiplicity.

Let $\mathcal N_g$ be the minimizing-geodesic nerve
constructed in Section~\ref{sec:nerve}, and let
$Q_g=|\mathcal N_g|/G$.
By Theorem~\ref{thm:genericity}, $\dim\mathcal N_g\leq N$.
The equivariant map $\Phi_g$ from
Proposition~\ref{prop:Phi} descends to
$\bar\Phi_g:B\to Q_g$.

For each $b=(x,y)\in C_g$, choose lifts
$(u,v)$ and an ordered minimizing-label set
\[
 S_b=\{a_0,\ldots,a_N\}.
\]
The associated top-dimensional simplex of
$\mathcal N_g$ is denoted $\sigma_b$.
Write
\[
 F_b=(L_{a_1}-L_{a_0},\ldots,L_{a_N}-L_{a_0}).
\]
The regularity condition says that
$dF_b:T_bB\to\R^N$ is invertible.
Define
\begin{equation}\label{eq:local-sign}
 \epsilon_b=\operatorname{sgn}\det dF_b\in\{+1,-1\}.
\end{equation}

\begin{lemma}\label{lem:local-degree}
The local degree of $\bar\Phi_g$ at $b$,
with target orientation determined by the
ordered simplex $[a_0,\ldots,a_N]$, equals
$\epsilon_b$.
\end{lemma}

\begin{proof}
Choose a small neighborhood of $b$ and compatible
lifts so that the only labels with
$L_a<m+\delta$ are the labels in $S_b$.
This is possible after shrinking the neighborhood
and the uniform allowance $\delta$, without
changing the nerve.

On this neighborhood the barycentric weights are
\[
 \lambda_i=\frac{m+\delta-L_{a_i}}{
       \sum_{j=0}^{N}(m+\delta-L_{a_j})},
 \qquad 0\leq i\leq N.
\]
All weights are strictly positive near $b$.
Their differences satisfy the exact identity
\begin{equation}\label{eq:barycentric-difference}
 \lambda_i-\lambda_0
 =-\frac{L_{a_i}-L_{a_0}}{W},
 \qquad 1\leq i\leq N.
\end{equation}
The barycenter is characterized by the vanishing
of the $N$ differences in
\eqref{eq:barycentric-difference}, so its inverse
image near $b$ is precisely $\{b\}$.
At $b$, $W=(N+1)\delta>0$. Although
$m=\min_a L_a$, and hence $W$, need not be
smooth at $b$, the continuous positive factor
$1/W$ does not affect the local Brouwer degree
of the map in \eqref{eq:barycentric-difference}.
Indeed, in local coordinates supplied by the
inverse function theorem for $F_b$, the
difference-coordinate map is
$t\mapsto-t/W(F_b^{-1}(t))$.
On a sufficiently small ball about the origin,
this is homotopic through maps with no boundary
zeros to $t\mapsto-t/W(b)$, by interpolating the
positive denominator with the constant $W(b)$.

The orientation defined by
$(\lambda_1-\lambda_0,\ldots,\lambda_N-\lambda_0)$
agrees with the standard simplex orientation:
the transition matrix from
$(\lambda_1,\ldots,\lambda_N)$ has
positive determinant $N+1$.
Therefore the local degree equals
$(-1)^N\operatorname{sgn}\det dF_b
=\operatorname{sgn}\det dF_b$, since $N=2n$.
No smoothness of $\Phi_g$ across the
minimum-length strata is required.
\end{proof}

\begin{lemma}\label{lem:localized-cycle}
The cellular cycle
\begin{equation}\label{eq:nerve-cycle}
 Z_g=\sum_{b\in C_g}\epsilon_b[\sigma_b]
 \in C_N(Q_g;\Z)
\end{equation}
represents $(\bar\Phi_g)_*[B]$ in $H_N(Q_g;\Z)$.
If multiple points yield the same orbit of
oriented simplices, their coefficients are added.
\end{lemma}

\begin{proof}
By Lemma~\ref{lem:no-inversions}, each open top
cell of $Q_g$ has an ordinary characteristic
simplex without vertex permutations. For each open top cell choose its barycenter.
Its inverse images are isolated points in $C_g$,
and Lemma~\ref{lem:local-degree} computes their
Brouwer local degrees. The cellular coefficient
of the pushforward of the fundamental class
is the sum of these degrees, by the relative
homology definition of degree. This reasoning
uses continuous maps and does not require
the barycenter to be a smooth regular value.
The resulting chain is \eqref{eq:nerve-cycle}.
It is a cycle because $B$ has no boundary.
\end{proof}

We now assemble the local contributions into a global formula.

\begin{theorem}\label{thm:localization}
Let $M$ be connected, closed, oriented and negatively curved,
and let $g$ be regular. Then
\begin{equation}\label{eq:localization}
 \mathfrak a_N\frown[M\times M]
 =\sum_{b\in C_g}\epsilon_b
 \big[(a_1-a_0)\wedge\cdots\wedge(a_N-a_0)\big]
 \quad\text{in }(\bigwedge^N I)_G.
\end{equation}
The right-hand side does not depend on the
chosen lifts or vertex orderings, provided
$\epsilon_b$ is defined by \eqref{eq:local-sign}.
\end{theorem}

\begin{proof}
Compose $\bar\Phi_g:B\to Q_g$ with the
simplicial inclusion $Q_g\to Q$, where
$Q=|Y|/G$ is the full simplex orbit space.
Proposition~\ref{prop:comparison} gives
\[
 \iota_N(\bar\Phi_g)_*[B]
 =\mathfrak a_N\frown[B].
\]
By Lemma~\ref{lem:localized-cycle}, the class
on the left is represented by $Z_g$.
Lemma~\ref{lem:exterior-quotient} sends
the simplex $[a_0,\ldots,a_N]$ to
\[
 (a_1-a_0)\wedge\cdots\wedge(a_N-a_0),
\]
which proves the formula.

Changing lifts applies the same $G$-element
to all labels, hence does not alter their
coinvariant class. Permuting the labels
changes the orientation of the simplex and
the sign of the determinant by the same
sign. The product appearing in
\eqref{eq:localization} is therefore unchanged.
\end{proof}

\begin{proposition}\label{prop:endpoint-reversal}
Let $J:\bigl(\bigwedge^N I\bigr)_G\to
\bigl(\bigwedge^N I\bigr)_G$ be the involution
induced by the additive map
\[
 \Z[\Gamma]\longrightarrow\Z[\Gamma],
 \qquad a\longmapsto a^{-1}.
\]
For $b=(x,y)\in C_g$, write
\[
 \tau_b=\epsilon_b
 \big[(a_1-a_0)\wedge\cdots\wedge(a_N-a_0)\big].
\]
Then the reversed pair $b^{\mathrm{op}}=(y,x)$
also belongs to $C_g$, and
\begin{equation}\label{eq:endpoint-reversal}
 \tau_{b^{\mathrm{op}}}=(-1)^n J(\tau_b).
\end{equation}
In particular $\#C_g$ is even and
\[
 J(A_M)=(-1)^n A_M.
\]
More precisely, $A_M$ lies in the image of
$1+(-1)^n J$ on the integral coinvariant module.
\end{proposition}

\begin{proof}
The inversion map sends the action
$(g,h)\cdot a=gah^{-1}$ to the action with
the two group coordinates interchanged:
\[
 (gah^{-1})^{-1}=h a^{-1}g^{-1}.
\]
It therefore induces the stated involution
on coinvariants and their exterior powers.

Choose lifts $u,v$ of $x,y$, and label the
minimizing geodesics by
$a_0,\ldots,a_N$. For the reversed endpoints
choose lifts $v,u$ and labels
$a_0^{-1},\ldots,a_N^{-1}$. The equality
\[
 d_X(v,a_i^{-1}u)=d_X(u,a_iv)
\]
shows that the reversed length-difference
map is the original length-difference map
composed with the endpoint swap
$T(u,v)=(v,u)$. On the $N=2n$ dimensional
product, $T$ has orientation degree
$(-1)^{n^2}=(-1)^n$. Hence the two
determinant signs satisfy
\[
 \epsilon_{b^{\mathrm{op}}}
 =(-1)^n\epsilon_b.
\]
Combining this with inversion of the
exterior generator gives
\eqref{eq:endpoint-reversal}.

The diagonal consists only of pairs joined
by one minimizing geodesic, so reversal acts
freely on $C_g$. The cardinality is even.
Summing \eqref{eq:endpoint-reversal}
over all points and using the localization
formula proves the eigenspace identity.
Selecting one point from each reversal
orbit also expresses $A_M$ as
$(1+(-1)^nJ)$ applied to an integral element.
\end{proof}

The local exterior terms admit a description in magnitude homology.

At each $b=(x,y)\in C_g$ the label set has
$N+1$ elements. Proposition
\ref{prop:MH-minimal} gives
\[
 \MH_{2,d_M(x,y)}(M;x,y)\cong I(S_b)
\]
with $\operatorname{rank}I(S_b)=N$.
The inclusion $I(S_b)\hookrightarrow I$
takes the generator
\[
 (a_1-a_0)\wedge\cdots\wedge(a_N-a_0)
 \in\bigwedge^N I(S_b)
\]
to the local exterior term in
\eqref{eq:localization}.
This interpretation is canonical up to the
orientation sign already accounted for
by $\epsilon_b$. Thus each summand is a
local index valued in the top exterior
power of second magnitude homology at
the minimal endpoint distance.

Combining localization with the seminorm estimate proves the counting theorem.

\begin{proof}[Proof of Theorem~\ref{thm:intro-count}]
Theorem~\ref{thm:genericity} proves that
regular metrics form an open dense set
and that $C_g$ is finite for every such metric.
By \eqref{eq:localization}, the real class
$A_M$ has a representation as a sum of
$\#C_g$ exterior generators, each with
coefficient $+1$ or $-1$. Definition
\eqref{eq:delta-norm} therefore yields
\[
 \|A_M\|_\Delta\leq\#C_g.
\]
The lower estimate from
Theorem~\ref{thm:intro-estimate} gives
\[
 \#C_g\geq\|A_M\|_\Delta
 \geq c_n\|M\|^2.
\]
The claim follows.
\end{proof}

\begin{corollary}\label{cor:top-multiplicity}
For every regular negatively curved metric on an
oriented closed negatively curved $n$-manifold,
there exists an ordered pair of points joined
by exactly $2n+1$ minimizing geodesics.
\end{corollary}

\begin{proof}
The simplicial volume of a closed negatively
curved manifold is positive, so the lower
bound in Theorem~\ref{thm:intro-count}
forces $\#C_g>0$.
\end{proof}

\begin{example}\label{ex:surface-count}
For a closed oriented surface $\Sigma_h$ of genus
$h\geq2$, one has $\|\Sigma_h\|=4(h-1)$.
Here $n=2$ and $c_2=1/72$. Therefore
every regular negatively curved metric satisfies
\[
 \#\{(x,y):m_g(x,y)=5\}
 \geq\frac{2}{9}(h-1)^2.
\]
The count concerns ordered pairs.
\end{example}

\begin{remark}
The estimate does not state that every negatively
curved metric is regular. At a nongeneric metric,
several top-dimensional local contributions may
coalesce at a point of higher multiplicity.
A general local-index formula for such degeneracies
requires additional analysis and is not needed
for the counting theorem proved here.
\end{remark}
